% lem-cell: limit-shapes.tex:466-471 (label lem:cell)
\paragraph{Paper excerpt.}
\begin{verbatim}
\begin{lemma}\label{lem:cell}
There exists $C(d)<\infty$ such that, for every $x\in\Z^d$,
\begin{equation*}
\int_{C_x}|\nabla \gauge(v)-\xi(x)|\dd v\leq C\mu\bigl(B(x,6\sqrt d)\bigr)\, .
\end{equation*}
\end{lemma}
\end{verbatim}
\paragraph{Lean statement.}
\begin{verbatim}
theorem CERW.Frozen.cell_gradient {d : ℕ} (hd : 2 ≤ d) :
    ∃ C : ℝ, 0 < C ∧
      ∀ Ψ : EuclideanSpace ℝ (Fin d) → ℝ, CERW.IsNorm Ψ →
      ∀ ξ : Site d → EuclideanSpace ℝ (Fin d),
        (∀ x : Site d, x ≠ 0 → CERW.IsSubgradient Ψ (CERW.toSpace x) (ξ x)) → ξ 0 = 0 →
      ∀ m : Measure (EuclideanSpace ℝ (Fin d)), IsLocallyFiniteMeasure m →
        CERW.IsDistribLaplacian Ψ m →
      ∀ x : Site d,
        ENNReal.ofReal (∫ v in CERW.cell x, ‖gradient Ψ v - ξ x‖)
          ≤ ENNReal.ofReal C * m (Metric.ball (CERW.toSpace x) (6 * Real.sqrt d))
--
\end{verbatim}
\paragraph{Transcription.}
Paper macros: \verb!\gauge! is $\Psi$. \emph{Object mapping.} \begin{itemize} \item $\gauge$, $\xi$: \verb!Ψ! with \verb!CERW.IsNorm Ψ!, and \verb!ξ : Site d → EuclideanSpace ℝ (Fin d)! with \verb!∀ x, x ≠ 0 → CERW.IsSubgradient Ψ (CERW.toSpace x) (ξ x)! and \verb!ξ 0 = 0! (line 303). The hypothesis \verb!ξ 0 = 0! is essential here: the lemma is asserted for every $x\in\Z^d$ including $x=0$, where $\xi(0)=0$ is a subgradient because $\gauge\ge0$ (paper proof, line 462). \verb!hd : 2 ≤ d! is the standing assumption. \item $C_x$ is \verb!CERW.cell x! $=\{v:\ \forall i,\ x_i-\tfrac12\le v_i<x_i+\tfrac12\}$ (half-open cube; the half-open or closed choice is immaterial for the Lebesgue integral). $B(x,6\sqrt d)$ is the open ball \verb!Metric.ball (CERW.toSpace x) (6 * Real.sqrt d)!. \item $\nabla\gauge(v)$ is \verb!gradient Ψ v!, the Mathlib gradient (value $0$ where $\gauge$ is not
